%
% Appendix: Linear Algebra
%

\chapterimage{Escher.pdf} % Chapter heading image

\chapter{Linear Algebra}

\begin{quote}
    \begin{flushright}
        \emph{All great work is the fruit of patience and perseverance,\\
            combined with tenacious concentration on a subject\\
            over a period of months or years.}\\
        Santiago Ramón y Cajal
    \end{flushright}
\end{quote}
\bigskip

{\color{red} TODO: Pending}

%
% Section: Vectors
%
\section{Vectors}


%
% Section: Matrices
%

\section{Matrices}

A \emph{matrix} $A$ of order $m \times n$ is a set of $mn$ scalars ordered in $m$ \emph{rows} and $n$ \emph{colums} in the following way:
\[
A = 
 \begin{pmatrix}
  a_{1,1} & a_{1,2} & \cdots & a_{1,n} \\
  a_{2,1} & a_{2,2} & \cdots & a_{2,n} \\
  \vdots  & \vdots  & \ddots & \vdots  \\
  a_{m,1} & a_{m,2} & \cdots & a_{m,n} 
 \end{pmatrix}
\]
The \emph{entry} $a_{ij}$ corresponds to the element of $A$ localed at row $i$ and column $j$. We denote by $\mathcal{M}_{m \times n}$ the \emph{set of all matrices} of order $m \times n$. A \emph{row matrix} is a matrix that belongs to $\mathcal{M}_{1 \times n}$, and a \emph{column matrix} to $\mathcal{M}_{m \times 1}$. A \emph{squared matrix} is a matrix that belongs to $\mathcal{M}_{n \times n}$. The entries $a_{ii}$ form the \emph{main diagonal} of a square matrix. The \emph{transpose matrix} of a matrix $A \in \mathcal{M}_{m \times n}$ is the matrix $A^T \in \mathcal{M}_{n \times m}$ whose entries $(i,j)$ are equal to the entries $(j, i)$ of $A$. If $A = A^T$ we say that $A$ is a \emph{symmetric matrix}. A \emph{diagonal matrix} is a matrix whose all its entries outside the main diagonal are zero. The \emph{identity matrix} $I$ is a diagonal squared matrix with all the entries in the main diagonal equal to 1. A \emph{submatrix} of a matrix is obtained by deleting any collection of rows and/or columns.

\begin{example}
Given the squared matrix $A = \left( \begin{smallmatrix} 1 & 2 & 3 \\ 4 & 5 & 6 \\ 7 & 8 & 9 \end{smallmatrix} \right)$ we have that the entry $(2, 3)$ has the value $6$, the diagonal is composed by the numbers $1$, $5$, and $9$, its transpose is $A^T = \left( \begin{smallmatrix} 1 & 4 & 7 \\ 2 & 5 & 8 \\ 3 & 6 & 9 \end{smallmatrix} \right)$, and that the matrix $B = \left( \begin{smallmatrix} 1 & 3 \\ 4 & 6 \end{smallmatrix} \right)$ is a submatrix of $A$.
\end{example}

The sum of two matrices $A$ and $B$ of equal size is the matrix $A + B$ whose entry $(i, j)$ are $(A + B)_{ij} = a_{ij} + b_{ij}$. The operation of sum of matrices is associative $(A + B) + C = A + (B + C)$, commutative $A + B = B + A$, has a neutral element $A + 0_{m \times n} = A$ and an inverse element $A + (-A) = 0_{m \times n}$. The multiplication of a scalar $\lambda$ by a matrix $A$ is the matrix $\lambda A$ whose entry $(i, j)$ is $(\lambda A)_{ij} = \lambda a_{ij}$. The operation of multiplication of a scalar by a matrix is distributive with respect to the sum of matrices $\alpha (A + B) = \alpha A + \alpha B$, distributive with respect to the sum of scalars $(\alpha + \beta) A = \alpha A + \beta B$, associative with respect to the product by scalars $(\alpha \beta) A = \alpha (\beta A)$, and has a unit element $1 A = A$. The product of two matrices $A_{m \times n}$ and $B_{n \times p}$ is the matrix $AB_{m \times p}$ whose entry $(i, j)$ is $(AB)_{ij} = \sum_{k=1}^n a_{ik} b_{kj}$. The operation of multiplication of matrices is associative $(A B) D = A (B D)$, has a left neutral element $A I_n = A$ and a right neutral element $I_m A = A$, associative with respect to the product by scalars $\alpha (A B) = (\alpha A) B = A (\alpha B)$, distributive with respect the sum of matrices by the right $A (B + C) = AB + AC$ and the left $(B + C) D = B D + C D$. Additionally, the operation of traspose satisfies the following properties: $(A + B)^T = A^T + B^T$, $(\lambda A)^T = \lambda A^T$, $(A B)^T = B^T A^T$.

\begin{example}
Given the matrices $A = \left( \begin{smallmatrix} 1 & 2 \\ 3 & 5 \end{smallmatrix} \right)$ and $B = \left( \begin{smallmatrix} 5 & 6 \\ 7 & 8 \end{smallmatrix} \right)$ we have that $A + B = \left( \begin{smallmatrix} 6 & 8 \\ 10 & 13 \end{smallmatrix} \right)$, $2 A = \left( \begin{smallmatrix} 2 & 4 \\ 6 & 10 \end{smallmatrix} \right)$ and $A B = \left( \begin{smallmatrix} 19 & 22 \\ 50 & 58 \end{smallmatrix} \right)$.
\end{example}

A square matrix $A$ is \emph{invertible} or \emph{non-singular} if there exists a matrix $B$ such that $AB = BA = I$. If $A$ is non-singular, $B$ is unique and is called the \emph{inverse matrix} of $A$, denoted by $A^{-1}$. A matrix $A$ is \emph{orthogonal} if its transpose is equal to its inverse $A^T = A^{-1}$; the columns and rows of a orthogonal matrix are called \emph{orthonormal vectors}.

The determinant is a scalar value that is a function of the entries of a square matrix denoted det(A)
\[
Leibniz formula.
\]
The determinant is nonzero if and only if the matrix is invertible

{\color{red} rank}

A number $\lambda$, and a non-zero vector $\mathbf{v}$ satisfying $A \mathbf{v} = \lambda \mathbf{v}$ are called an \emph{eigenvalue} and an \emph{eigenvector} of $A$, respectively.

\begin{example}
The determinant of a 2 × 2 matrix is
% \begin{vmatrix}a&b\\c&d\end{vmatrix}}=ad-bc
and the determinant of a 3 × 3 matrix is
% \begin{vmatrix}a&b&c\\d&e&f\\g&h&i\end{vmatrix}}=aei+bfg+cdh-ceg-bdi-afh.
Refer to the References section 
\end{example}

Matrix decomposition is the process to render a matrix into a more easily accessible form, meanwhile certain properties, like the determinat or the rank, are preserved. The \emph{singular value decomposition} of a matrix $A$ of order $m \times n$ is a factorization of the form $A = U \Sigma V^T$ where $U$ is an $m \times m$ orthogonal matrix, $\Sigma$ is an $m \times n$ non-negative rectangular diagonal matrix, and $V^T$ is the traspose of an $n \times n$ orthogogal matrix. 



%
% Section: References
%

\section*{References}

{\color{red} TODO: Pending.}
